\documentclass{beamer}

% Theme choice
\usetheme{Madrid} % You can change to e.g., Warsaw, Berlin, CambridgeUS, etc.

% Encoding and font
\usepackage[utf8]{inputenc}
\usepackage[T1]{fontenc}

% Graphics and tables
\usepackage{graphicx}
\usepackage{booktabs}

% Code listings
\usepackage{listings}
\lstset{
basicstyle=\ttfamily\small,
keywordstyle=\color{blue},
commentstyle=\color{gray},
stringstyle=\color{red},
breaklines=true,
frame=single
}

% Math packages
\usepackage{amsmath}
\usepackage{amssymb}

% Colors
\usepackage{xcolor}

% TikZ and PGFPlots
\usepackage{tikz}
\usepackage{pgfplots}
\pgfplotsset{compat=1.18}
\usetikzlibrary{positioning}

% Hyperlinks
\usepackage{hyperref}

% Title information
\title{Chapter 10: Debugging and Testing}
\author{Your Name}
\institute{Your Institution}
\date{\today}

\begin{document}

\frame{\titlepage}

\begin{frame}[fragile]
    \frametitle{Introduction to Debugging and Testing}
    \begin{block}{Understanding Debugging and Testing}
        \begin{itemize}
            \item \textbf{Definition of Debugging:} 
                Debugging is the process of identifying, isolating, and fixing problems ("bugs") in code.
            \item \textbf{Definition of Testing:} 
                Testing is the systematic process of executing a program to identify bugs or verify requirements.
        \end{itemize}
    \end{block}
\end{frame}

\begin{frame}[fragile]
    \frametitle{Importance in Programming}
    \begin{enumerate}
        \item \textbf{Code Quality:}
        \begin{itemize}
            \item \textbf{Correctness:} Ensures software meets its requirements.
            \item \textbf{Reliability:} Decreases the likelihood of failures.
            \item \textbf{Maintainability:} Simplifies future modifications.
        \end{itemize}
        
        \item \textbf{Performance:}
        \begin{itemize}
            \item \textbf{Responsiveness:} Faster applications improve user satisfaction.
            \item \textbf{Scalability:} Well-tested code can handle increased loads.
        \end{itemize}
    \end{enumerate}
\end{frame}

\begin{frame}[fragile]
    \frametitle{Key Points to Emphasize}
    \begin{itemize}
        \item \textbf{Proactive vs. Reactive:} 
            Effective debugging and testing prevent issues before reaching production.
        \item \textbf{The Development Cycle:} 
            Should be integrated throughout the entire software development lifecycle (SDLC).
        \item \textbf{Diverse Techniques:} 
            Techniques vary; unit tests for small segments, integration tests for module interactions.
    \end{itemize}
\end{frame}

\begin{frame}[fragile]
    \frametitle{Examples of Debugging and Testing}
    \begin{block}{Example of Debugging:}
        \begin{lstlisting}[language=python]
def calculate_area(radius):
    return 3.14 * radius ** 2

# Debugging example with a wrong expected call
print(calculate_area('five'))  # This will raise a TypeError
        \end{lstlisting}
        Debugging would involve ensuring a valid numeric input.
    \end{block}
\end{frame}

\begin{frame}[fragile]
    \frametitle{Examples of Debugging and Testing (Continued)}
    \begin{block}{Example of Testing:}
        \begin{lstlisting}[language=python]
import unittest

class TestCalculator(unittest.TestCase):
    def test_area(self):
        self.assertAlmostEqual(calculate_area(5), 78.5)

# Running this test ensures that calculate_area works as expected.
if __name__ == '__main__':
    unittest.main()
        \end{lstlisting}
    \end{block}
\end{frame}

\begin{frame}[fragile]
    \frametitle{Conclusion}
    \begin{block}{Conclusion:}
        Debugging and testing are essential for producing a robust, high-quality product. 
        Mastering these practices enhances both the product's success and the developer's efficiency.
    \end{block}
\end{frame}

\begin{frame}[fragile]
    \frametitle{Common Debugging Techniques - Overview}
    \begin{itemize}
        \item Debugging is essential for identifying and fixing errors in code.
        \item Discussing three common techniques:
        \begin{itemize}
            \item Print Statements
            \item Using Debuggers
            \item Logging
        \end{itemize}
        \item Each method has strengths and ideal use cases.
    \end{itemize}
\end{frame}

\begin{frame}[fragile]
    \frametitle{Common Debugging Techniques - Print Statements}
    \begin{block}{Description}
        The simplest and oldest debugging technique, involving print statements to output variable values and track execution flow.
    \end{block}
    
    \begin{block}{Example}
        \begin{lstlisting}[language=Python]
def add_numbers(a, b):
    print(f"Adding {a} and {b}")
    return a + b

result = add_numbers(5, 3)
print(f"The result is {result}")
        \end{lstlisting}
    \end{block}

    \begin{itemize}
        \item \textbf{Pros:}
        \begin{itemize}
            \item Fast and straightforward.
            \item No special tools required.
        \end{itemize}
        \item \textbf{Cons:}
        \begin{itemize}
            \item Can clutter code if overused.
            \item Not suitable for complex debugging.
        \end{itemize}
    \end{itemize}
\end{frame}

\begin{frame}[fragile]
    \frametitle{Common Debugging Techniques - Using Debuggers}
    \begin{block}{Description}
        Debuggers are tools that allow step-by-step execution, variable inspection, and control of logic flow.
    \end{block}
    
    \begin{block}{Features}
        \begin{itemize}
            \item Breakpoints: Pause execution at certain lines.
            \item Step-Over: Execute the next line without diving into functions.
            \item Watch Variables: Monitor values of selected variables.
        \end{itemize}
    \end{block}

    \begin{block}{Example (Using a Debugger)}
        Setting a breakpoint in Visual Studio Code to pause execution and inspect variables.
    \end{block}

    \begin{itemize}
        \item \textbf{Pros:}
        \begin{itemize}
            \item Powerful for complex applications.
            \item Visual and interactive to inspect execution.
        \end{itemize}
        \item \textbf{Cons:}
        \begin{itemize}
            \item Requires tool familiarity.
            \item Setup varies by language and IDE.
        \end{itemize}
    \end{itemize}
\end{frame}

\begin{frame}[fragile]
    \frametitle{Common Debugging Techniques - Logging}
    \begin{block}{Description}
        Logging is the recording of events during program execution, providing different levels (info, warning, error).
    \end{block}
    
    \begin{block}{Example}
        \begin{lstlisting}[language=Python]
import logging

logging.basicConfig(level=logging.INFO)

def multiply_numbers(a, b):
    logging.info(f"Multiplying {a} by {b}")
    return a * b

result = multiply_numbers(4, 2)
logging.info(f"The result is {result}")
        \end{lstlisting}
    \end{block}

    \begin{itemize}
        \item \textbf{Pros:}
        \begin{itemize}
            \item Provides a permanent record.
            \item Can filter messages by severity.
        \end{itemize}
        \item \textbf{Cons:}
        \begin{itemize}
            \item Requires logging framework setup.
            \item Can introduce performance overhead if unmanaged.
        \end{itemize}
    \end{itemize}
\end{frame}

\begin{frame}[fragile]
    \frametitle{Summary of Key Points}
    \begin{itemize}
        \item Different debugging techniques serve various purposes in software development.
        \item Print statements are suitable for quick checks.
        \item Debuggers and logging help manage complex debugging tasks.
        \item Mastering these techniques enhances error identification and code quality.
    \end{itemize}
\end{frame}

\begin{frame}
    \frametitle{Introduction to Testing}
    \begin{block}{Overview of the Role of Testing in Software Development}
        Testing is a critical phase in the software development lifecycle (SDLC) that involves evaluating software applications to ensure they perform as expected. The primary roles of testing include:
    \end{block}
    \begin{enumerate}
        \item \textbf{Error Detection}: Identifying bugs and errors in code before deployment or integration.
        \item \textbf{Quality Assurance}: Ensuring that the software meets specified requirements and standards.
        \item \textbf{Risk Management}: Mitigating potential risks associated with software failures which can lead to financial loss, security breaches, or user dissatisfaction.
    \end{enumerate}
\end{frame}

\begin{frame}
    \frametitle{Types of Errors that Testing Can Mitigate - Part 1}
    \begin{block}{Various Types of Errors}
        Testing can help identify a variety of errors, which can broadly be categorized as follows:
    \end{block}
    \begin{enumerate}
        \item \textbf{Syntax Errors}: Mistakes in the code that violate the grammatical rules of the programming language. 
        \begin{itemize}
            \item \textit{Example}:
            \begin{lstlisting}[language=Python]
print("Hello World"  # SyntaxError: missing parentheses in call to 'print'
            \end{lstlisting}
        \end{itemize}
        
        \item \textbf{Logic Errors}: These occur when the program runs without crashing but produces incorrect results.
        \begin{itemize}
            \item \textit{Example}:
            \begin{lstlisting}[language=Python]
total = price * quantity  # If intended logic was to add prices, should be total = price + quantity
            \end{lstlisting}
        \end{itemize}
    \end{enumerate}
\end{frame}

\begin{frame}
    \frametitle{Types of Errors that Testing Can Mitigate - Part 2}
    \begin{enumerate}
        \setcounter{enumi}{2} % Continue from previous enumeration
        \item \textbf{Runtime Errors}: Errors that only occur during execution.
        \begin{itemize}
            \item \textit{Example}:
            \begin{lstlisting}[language=Python]
array = [1, 2, 3]
print(array[3])  # IndexError: list index out of range
            \end{lstlisting}
        \end{itemize}

        \item \textbf{Performance Errors}: Issues that may degrade performance (e.g., memory leaks).
        \begin{itemize}
            \item \textit{Example}: Inefficient algorithms that lead to slow response times.
        \end{itemize}
    \end{enumerate}
\end{frame}

\begin{frame}
    \frametitle{Key Points to Emphasize}
    \begin{itemize}
        \item \textbf{Early Testing}: Incorporate testing early in the SDLC.
        \item \textbf{Automated Testing}: Utilize automated tests for repetitive tasks.
        \item \textbf{Different Types of Testing}: Understanding testing methodologies enhances error mitigation.
    \end{itemize}
    \begin{block}{Conclusion}
        Testing is an integral part of the software development process that contributes to the overall success of the project.
    \end{block}
\end{frame}

\begin{frame}[fragile]
    \frametitle{Types of Testing - Overview}
    \begin{block}{Overview of Testing Methodologies}
        Testing is a critical phase in software development aimed at identifying and rectifying errors before deployment. The three key types of testing methodologies are:
        \begin{itemize}
            \item \textbf{Unit Testing}
            \item \textbf{Integration Testing}
            \item \textbf{System Testing}
        \end{itemize}
    \end{block}
\end{frame}

\begin{frame}[fragile]
    \frametitle{Types of Testing - Unit Testing}
    \begin{block}{Unit Testing}
        \textbf{Definition}: Testing individual components or modules of the software. \\
        \textbf{Purpose}: Validate each unit performs as expected to prevent early bugs. 
    \end{block}
    \begin{itemize}
        \item \textbf{Implementation:}
            \begin{itemize}
                \item Written by developers using frameworks (e.g., JUnit, pytest).
                \item Each test case checks specific inputs against expected outputs.
            \end{itemize}
    \end{itemize}
    \begin{block}{Example}
        \begin{lstlisting}[language=Python]
def add(a, b):
    return a + b

# Unit test for the add function
def test_add():
    assert add(2, 3) == 5
    assert add(-1, 1) == 0
        \end{lstlisting}
    \end{block}
\end{frame}

\begin{frame}[fragile]
    \frametitle{Types of Testing - Integration Testing}
    \begin{block}{Integration Testing}
        \textbf{Definition}: Testing interactions between multiple components. \\
        \textbf{Purpose}: Detect interface defects when units are combined.
    \end{block}
    \begin{itemize}
        \item \textbf{Implementation:}
            \begin{itemize}
                \item Simulates real-world scenarios (e.g., Postman for APIs).
                \item Tests groups of units collectively. 
            \end{itemize}
    \end{itemize}
    \begin{block}{Example}
        \begin{lstlisting}[language=Python]
def fetch_user_details(user_id):
    user = get_user_from_db(user_id)
    profile = get_user_profile(user_id)
    return {**user, **profile}

# Integration test for fetching user details
def test_fetch_user_details():
    # Mocking database and profile responses
    user_data = {'id': 1, 'name': 'Jane Doe'}
    profile_data = {'age': 30, 'location': 'New York'}
    assert fetch_user_details(1) == {**user_data, **profile_data}
        \end{lstlisting}
    \end{block}
\end{frame}

\begin{frame}[fragile]
    \frametitle{Types of Testing - System Testing}
    \begin{block}{System Testing}
        \textbf{Definition}: Evaluating the entire system's compliance with requirements. \\
        \textbf{Purpose}: Ensure the complete software product meets specifications.
    \end{block}
    \begin{itemize}
        \item \textbf{Implementation:}
            \begin{itemize}
                \item Conducted in a production-like environment.
                \item Tests user interfaces, APIs, databases, and overall performance. 
            \end{itemize}
    \end{itemize}
    \begin{block}{Example}
        Performing a stress test on a web application to ensure it can handle 1000 concurrent users without failures.
    \end{block}
\end{frame}

\begin{frame}[fragile]
    \frametitle{Types of Testing - Conclusion & Next Steps}
    \begin{block}{Conclusion}
        Understanding Unit Testing, Integration Testing, and System Testing is essential for delivering high-quality software. They collectively ensure robustness, functionality, and user satisfaction in the final product.
    \end{block}
    \begin{block}{Next Steps}
        In the upcoming slide, we will delve into \textbf{Test-Driven Development (TDD)}, exploring its principles and benefits in the software development lifecycle.
    \end{block}
\end{frame}

\begin{frame}
    \frametitle{Test-Driven Development (TDD)}
    \begin{block}{Understanding TDD}
        Test-Driven Development (TDD) is a software development methodology that emphasizes writing tests before implementing code, following a cycle known as Red-Green-Refactor.
    \end{block}
\end{frame}

\begin{frame}
    \frametitle{The TDD Cycle}
    \begin{enumerate}
        \item \textbf{Red}: Write a test for the new functionality that should initially fail.
        \item \textbf{Green}: Write the minimal code to pass the test.
        \item \textbf{Refactor}: Improve the code while keeping all tests passing.
    \end{enumerate}
\end{frame}

\begin{frame}[fragile]
    \frametitle{Example: TDD Cycle}
    \begin{block}{Red Phase Example}
        \begin{lstlisting}[language=python]
def test_add_numbers():
    assert add(1, 2) == 3  # Test fails initially
        \end{lstlisting}
    \end{block}
    
    \begin{block}{Green Phase Example}
        \begin{lstlisting}[language=python]
def add(a, b):
    return a + b  # Code to pass the test
        \end{lstlisting}
    \end{block}
    
    \begin{block}{Refactor Phase Example}
        % Add refactoring explanation or example here
        \begin{lstlisting}[language=python]
# Refactoring might look like this:
def add_numbers(numbers):
    return sum(numbers)  # Refactored implementation
        \end{lstlisting}
    \end{block}
\end{frame}

\begin{frame}
    \frametitle{Benefits and Principles of TDD}
    \begin{itemize}
        \item \textbf{Improved Code Quality}: Reduces the likelihood of bugs.
        \item \textbf{Better Design}: Encourages modular and understandable code.
        \item \textbf{Documentation}: Tests serve as living documentation.
        \item \textbf{Facilitates Refactoring}: Provides confidence to change code safely.
    \end{itemize}
\end{frame}

\begin{frame}
    \frametitle{Using Python Testing Frameworks}
    \begin{block}{Introduction}
        Python offers several powerful testing frameworks that simplify the process of writing and running tests to ensure that code functions as intended. Two of the most popular frameworks are \textbf{unittest} and \textbf{pytest}.
    \end{block}
\end{frame}

\begin{frame}[fragile]
    \frametitle{Unittest}
    \begin{itemize}
        \item \textbf{What is it?}
        \begin{itemize}
            \item A built-in Python module that follows the xUnit style framework for writing and running tests.
        \end{itemize}
        
        \item \textbf{Key Features}:
        \begin{itemize}
            \item Supports test discovery and organization.
            \item Allows for setup and teardown of test cases.
            \item Provides various assertion methods to verify outcomes.
        \end{itemize}

        \item \textbf{Basic Example}:
        \begin{lstlisting}
import unittest

def add(a, b):
    return a + b

class TestMathOperations(unittest.TestCase):
    def test_add(self):
        self.assertEqual(add(2, 3), 5)
        self.assertEqual(add(-1, 1), 0)

if __name__ == '__main__':
    unittest.main()
        \end{lstlisting}
    \end{itemize}
\end{frame}

\begin{frame}
    \frametitle{Explanation of Unittest Example}
    The code defines a simple function \texttt{add} and a unit test class \texttt{TestMathOperations}. The test case checks various scenarios to ensure the \texttt{add} function produces correct output.
\end{frame}

\begin{frame}[fragile]
    \frametitle{Pytest}
    \begin{itemize}
        \item \textbf{What is it?}
        \begin{itemize}
            \item An advanced testing framework that facilitates simple and scalable testing while allowing for plugins.
        \end{itemize}

        \item \textbf{Key Features}:
        \begin{itemize}
            \item Requires less boilerplate code compared to unittest.
            \item Supports fixtures for setup code and has a rich plugin architecture.
            \item Easy to assert using simple Python assertions (\texttt{assert}).
        \end{itemize}

        \item \textbf{Basic Example}:
        \begin{lstlisting}
def add(a, b):
    return a + b

def test_add():
    assert add(2, 3) == 5
    assert add(-1, 1) == 0
        \end{lstlisting}
    \end{itemize}
\end{frame}

\begin{frame}
    \frametitle{Explanation of Pytest Example}
    Here, the \texttt{test_add} function directly uses assertions to validate the \texttt{add} function, demonstrating how pytest enables clearer, more concise test writing.
\end{frame}

\begin{frame}
    \frametitle{Key Points to Emphasize}
    \begin{itemize}
        \item \textbf{Test-Driven Development (TDD)}: Using these frameworks is aligned with the TDD approach where tests are written before the actual code, helping identify and fix bugs early.
        \item \textbf{Suitability}: Choose \texttt{unittest} for more structured, traditional testing. Opt for \texttt{pytest} for flexibility, ease of use, and extensive testing capabilities.
        \item \textbf{Running Tests}: 
        \begin{itemize}
            \item For unittest: \texttt{python -m unittest discover}
            \item For pytest: simply run \texttt{pytest} in the terminal.
        \end{itemize}
    \end{itemize}
\end{frame}

\begin{frame}
    \frametitle{Conclusion}
    Understanding and utilizing testing frameworks like \textbf{unittest} and \textbf{pytest} is crucial for maintaining high-quality code. They facilitate not just the identification of bugs but also help in building a robust codebase through consistent testing practices.
\end{frame}

\begin{frame}[fragile]
    \frametitle{Debugging in Python}
    \begin{block}{Understanding Debugging}
        Debugging is the process of identifying and removing errors or bugs from the code.
        It is a critical skill for Python developers as it helps ensure code runs smoothly and as intended.
    \end{block}
\end{frame}

\begin{frame}[fragile]
    \frametitle{Setting Breakpoints}
    \begin{itemize}
        \item \textbf{What is a Breakpoint?}  
        A breakpoint is a designated stopping point in your code where the execution will pause, allowing 
        you to inspect the program's state (variables, memory, etc.) at that precise moment.
        
        \item \textbf{How to Set a Breakpoint in IDEs:}
        \begin{enumerate}
            \item \textbf{Visual Studio Code (VS Code):}
            \begin{enumerate}
                \item Open the file you want to debug.
                \item Click in the gutter next to the line where you want to set the breakpoint; a red dot will appear.
                \item Start debugging by pressing \texttt{F5} or selecting "Run" > "Start Debugging".
            \end{enumerate}
            \item \textbf{PyCharm:}
            \begin{enumerate}
                \item Open your Python file.
                \item Click in the left margin next to the line number where you want the code to stop.
                \item Begin debugging by clicking the bug icon in the toolbar or pressing \texttt{Shift + F9}.
            \end{enumerate}
        \end{enumerate}
    \end{itemize}
\end{frame}

\begin{frame}[fragile]
    \frametitle{Stepping Through Code}
    \begin{itemize}
        \item \textbf{Step Over (F10):} Executes the next line but doesn’t enter any called functions.
        
        \item \textbf{Step Into (F11):} Enters the next function call to see its inner workings.
        
        \item \textbf{Step Out (Shift + F11):} Runs until the current function exits and returns to the previous level.
        
        \item \textbf{Continue (F5):} Resumes execution until the next breakpoint is reached.
    \end{itemize}
\end{frame}

\begin{frame}[fragile]
    \frametitle{Example Code for Debugging}
    \begin{lstlisting}[language=Python]
def divide_numbers(num1, num2):
    # Set a breakpoint here
    return num1 / num2

result = divide_numbers(10, 0)  # This will raise an error
print(result)
    \end{lstlisting}

    \begin{block}{Concept}
        Setting a breakpoint on the line with \texttt{return num1 / num2} allows you to 
        inspect \texttt{num1} and \texttt{num2} before the division occurs and 
        catch the error raised by dividing by zero.
    \end{block}
\end{frame}

\begin{frame}[fragile]
    \frametitle{Key Points to Emphasize}
    \begin{itemize}
        \item Use of Breakpoints provides a powerful mechanism to pause execution to inspect variable states and control flow.
        
        \item Understanding the difference between Step Over, Step Into, and Step Out is crucial for effective debugging.
        
        \item Most modern IDEs like VS Code and PyCharm offer intuitive interfaces for setting breakpoints and debugging, which can greatly enhance productivity.
    \end{itemize}
\end{frame}

\begin{frame}[fragile]
    \frametitle{Conclusion}
    Effective debugging is essential for developing robust Python applications. 
    By mastering breakpoints and step-through debugging, you can ensure your applications run efficiently and handle errors gracefully. 

    \begin{block}{Encouragement}
        Encourage students to practice debugging through examples in their own code and explore the features of 
        their preferred IDE for a deeper understanding of these concepts.
    \end{block}
\end{frame}

\begin{frame}
    \frametitle{Best Practices for Debugging and Testing}
    \begin{block}{Effective Strategies for Reliable Code}
        \begin{enumerate}
            \item Understand Your Code
            \item Use a Debugger
            \item Isolate the Problem
            \item Test Early and Often
            \item Use Version Control
            \item Collaborate and Review
            \item Maintain a Debugging Mindset
        \end{enumerate}
    \end{block}
\end{frame}

\begin{frame}[fragile]
    \frametitle{Understanding Your Code}
    \begin{itemize}
        \item \textbf{Read Before You Debug:} Familiarize yourself with the codebase and logical flow.
        \item \textbf{Comment and Document:} Explain complex parts and document functions clearly.
    \end{itemize}
\end{frame}

\begin{frame}[fragile]
    \frametitle{Using a Debugger}
    \begin{itemize}
        \item \textbf{IDE Features:} Utilize built-in debugging tools such as breakpoints, step execution, and variable inspection.
        \item \textbf{Example:} In Python, tools like PyCharm or Visual Studio Code provide functionalities for setting breakpoints and evaluating expressions on the fly.
    \end{itemize}
\end{frame}

\begin{frame}[fragile]
    \frametitle{Isolating the Problem}
    \begin{itemize}
        \item \textbf{Divide and Conquer:} Break down code to identify the source of the bug.
        \item \textbf{Print Statements:} Track variable values and program flow using print statements or logging.
    \end{itemize}
\end{frame}

\begin{frame}[fragile]
    \frametitle{Testing Early and Often}
    \begin{itemize}
        \item \textbf{Unit Testing:} Write unit tests for functions to catch bugs early in the development cycle.
        \item \textbf{Example:}
        \begin{lstlisting}[language=Python]
import unittest

class TestSum(unittest.TestCase):
    def test_sum(self):
        self.assertEqual(sum([1, 2, 3]), 6)

if __name__ == '__main__':
    unittest.main()
        \end{lstlisting}
    \end{itemize}
\end{frame}

\begin{frame}
    \frametitle{Additional Best Practices}
    \begin{itemize}
        \item \textbf{Use Version Control:} Manage changes with Git or other systems to revert to previous states easily.
        \item \textbf{Collaborate and Review:} Engage in peer reviews and pair programming to catch overlooked errors.
        \item \textbf{Maintain a Debugging Mindset:} 
            \begin{itemize}
                \item Be patient and persistent; maintain a calm and methodical approach.
                \item Take breaks if stuck to find new perspectives.
            \end{itemize}
    \end{itemize}
\end{frame}

\begin{frame}
    \frametitle{Key Points to Emphasize}
    \begin{itemize}
        \item \textbf{Documentation is Vital:} Aids in debugging and future development.
        \item \textbf{Testing is Not Just a Phase:} Integrate it throughout the development lifecycle.
        \item \textbf{Learning from Errors:} Treat bugs as learning opportunities for skill improvement.
    \end{itemize}
\end{frame}

\begin{frame}[fragile]
    \frametitle{Common Debugging Challenges - Overview}
    \begin{itemize}
        \item Debugging is crucial in software development.
        \item Developers face various challenges that hinder effectiveness. 
        \item Understanding these challenges and strategies to overcome them is key.
    \end{itemize}
\end{frame}

\begin{frame}[fragile]
    \frametitle{Common Debugging Challenges - Part 1}
    \begin{enumerate}
        \item \textbf{Over-Complexity of Code}
            \begin{itemize}
                \item \textbf{Challenge:} Difficult to trace logic.
                \item \textbf{Solution:} Refactor code; simplify with clear naming.
                \item \textbf{Example:} 
                \begin{lstlisting}[language=Python]
def process(order):
    if order.status == "Pending":
        # lengthy conditions
        if order.amount > 1000 and order.country == "USA":
            pass

# Refactored
def handle_pending_order(order):
    if order.amount > 1000 and order.country == "USA":
        process_large_order(order)
                \end{lstlisting}
            \end{itemize}
        
        \item \textbf{Misleading Error Messages}
            \begin{itemize}
                \item \textbf{Challenge:} Vague error messages confuse diagnosis.
                \item \textbf{Solution:} Improve error logging with context.
                \item \textbf{Key Points:}
                    \begin{itemize}
                        \item Use descriptive error messages.
                        \item Log stack traces for better tracking.
                    \end{itemize}
            \end{itemize}
    \end{enumerate}
\end{frame}

\begin{frame}[fragile]
    \frametitle{Common Debugging Challenges - Part 2}
    \begin{enumerate}
        \setcounter{enumi}{2}
        \item \textbf{Intermittent Bugs}
            \begin{itemize}
                \item \textbf{Challenge:} Bugs that appear sporadically.
                \item \textbf{Solution:} Capture environmental conditions.
                \item \textbf{Example:} Implement logging for variable states.
            \end{itemize}

        \item \textbf{Failure to Isolate the Problem}
            \begin{itemize}
                \item \textbf{Challenge:} Bugs propagate through systems.
                \item \textbf{Solution:} Use unit tests and mock components.
                \item \textbf{Key Points:}
                    \begin{itemize}
                        \item Document tests and expected behaviors.
                    \end{itemize}
            \end{itemize}

        \item \textbf{Ignoring Edge Cases}
            \begin{itemize}
                \item \textbf{Challenge:} Overlooking unusual inputs.
                \item \textbf{Solution:} Employ boundary testing, validate inputs.
                \item \textbf{Example:} Consider scenarios for numeric inputs.
            \end{itemize}
    \end{enumerate}
\end{frame}

\begin{frame}[fragile]
    \frametitle{Common Debugging Challenges - Part 3}
    \begin{enumerate}
        \setcounter{enumi}{5}
        \item \textbf{Sticking to a Single Debugging Method}
            \begin{itemize}
                \item \textbf{Challenge:} Limiting effectiveness with one method.
                \item \textbf{Solution:} Utilize various debugging strategies.
                \item \textbf{Key Points:}
                    \begin{itemize}
                        \item Experiment with different debugging techniques.
                        \item Find the right tool for the problem.
                    \end{itemize}
            \end{itemize}
    \end{enumerate}
    
    \begin{block}{Conclusion}
        Recognizing these debugging challenges is crucial for a productive development environment. By applying strategic approaches, developers enhance their problem-solving skills and improve overall code quality.
    \end{block}
\end{frame}

\begin{frame}[fragile]
    \frametitle{Conclusion and Summary - Key Points Recap}
    \begin{enumerate}
        \item \textbf{Importance of Debugging and Testing}
        \begin{itemize}
            \item Debugging identifies and resolves defects, crucial for reliable software.
            \item Testing ensures software meets requirements and maintains quality.
        \end{itemize}
        \item \textbf{Stages of Debugging}
        \begin{itemize}
            \item Reproduction of the Issue
            \item Diagnosis
            \item Resolution
            \item Verification
        \end{itemize}
        \item \textbf{Types of Testing}
        \begin{itemize}
            \item Unit Testing
            \item Integration Testing
            \item System Testing
            \item User Acceptance Testing (UAT)
        \end{itemize}
    \end{enumerate}
\end{frame}

\begin{frame}[fragile]
    \frametitle{Conclusion and Summary - Tools, Techniques, and Best Practices}
    \begin{enumerate}
        \setcounter{enumi}{3}
        \item \textbf{Tools and Techniques}
        \begin{itemize}
            \item Use testing frameworks (e.g., JUnit, pytest) for automation.
            \item Implement CI/CD practices for efficiency.
        \end{itemize}
        \item \textbf{Best Practices}
        \begin{itemize}
            \item Write tests first (Test-Driven Development - TDD).
            \item Regularly review and refactor code.
            \item Document debugging findings.
        \end{itemize}
    \end{enumerate}
\end{frame}

\begin{frame}[fragile]
    \frametitle{Conclusion and Summary - Significance in Software Development}
    \begin{block}{Significance in the Software Development Lifecycle}
        \begin{itemize}
            \item \textbf{Risk Mitigation:} Reduces software failure risks and builds user trust.
            \item \textbf{Cost Efficiency:} Early debugging lowers maintenance costs.
            \item \textbf{Enhanced Collaboration:} Clear tested codebase facilitates teamwork.
        \end{itemize}
    \end{block}

    \begin{block}{Conclusion}
        Mastering debugging and testing is essential for quality software delivery, user satisfaction, and cost reduction.
    \end{block}
\end{frame}


\end{document}